\subsubsection{A necessary condition for a finite cone cover}

\begin{proposition}[Proved: A persistent first neighborhood: proved obstruction]
\label{D:persistent-first-neighborhood}
Let $k$ be algebraically closed of characteristic $p>0$, and let
$C$ be a smooth connected projective curve of genus $g\geq2$ in a smooth
surface $Y$.  Put $N=N_{C/Y}$ and $d=-\deg N$, and assume
\[
                         0<d<2g-2.
\]
Let
\[
 e_Y\in H^1(C,T_C\otimes N^{-1})
\]
be the class of the first infinitesimal neighborhood of $C$.
Suppose that every Frobenius pullback of $e_Y$ is nonzero.
Then there is no finite surjective adic morphism
\[
 \phi:\widehat{\operatorname{Tot}(N')}_{C'}
             \longrightarrow\widehat Y_C
\]
where $C'$ is a smooth connected projective curve and $N'$ is a line
bundle on $C'$.  Transverse ramification is allowed in this assertion.

In particular, the hypothesis on Frobenius is satisfied whenever
\[
                   e_Y\ne0,\qquad p(2g-2-d)>2g-2.
\]
This is an obstruction to the finite-cone method.  We make no assertion
here about finite generation of the corresponding punctured fundamental
group.
\end{proposition}

\begin{proof}
We first record the curve-theoretic input.  If $L$ is a negative-degree
line bundle on a smooth projective curve and
$c\in H^1(C,L)$ survives all Frobenius pullbacks, its additive $L$-torsor
is affine.  This is the rank-one case of Brenner's affineness criterion,
valid over an arbitrary algebraically closed field of positive
characteristic: \cite[Proposition 2.1, preprint pp.~5--6]{C_Brenner2006}.
Consequently $h^*c\ne0$ for every finite surjective morphism
$h:C'\to C$ from a smooth projective connected curve.  Otherwise a
section of the pulled-back torsor would give a nonconstant morphism
from the proper curve $C'$ to the affine torsor, which is impossible.
Conversely, if a Frobenius power kills $c$, that power is itself a
finite surjective curve morphism.  Thus for a negative line, finite
annihilation and Frobenius annihilation are equivalent.  We use only
the implication just proved, without a finite-field hypothesis.

Suppose that $\phi$ exists, and let $h:C'\to C$ be its restriction to
the reduced closed fibers.  Finiteness and surjectivity make $h$ finite
surjective; write $m=\deg h$.  Because $\phi$ is adic, the pullback of
the ideal of $C$ has the same radical as the ideal of the zero section
$C'$.  Both formal surfaces are regular.  The divisor of a pulled-back
local equation of $C$ is therefore
\[
                         \phi^*C=eC'
\]
for one positive integer $e$.  In particular its leading coefficient
is a unit, and restriction of the associated invertible ideals gives
\[
                    (N')^{\otimes e}\simeq h^*N,
             \qquad \deg N'=-\frac{md}{e}.             \tag{PN1}
\]
This is the point at which the adic and finite hypotheses exclude an
additional horizontal component in the inverse image of $C$.

For $n\geq1$, write $C'_n$ for the zero section thickened by the
$n$th power of its ideal.  Since a local equation of $C$ maps into
the $e$th power of that ideal, the restriction of $\phi$ to $C'_e$
factors through $C$.  We claim that this factorization is the composite
of the canonical cone retraction $C'_e\to C'$ and $h$.
Both maps agree on $C'$.  At the step from $C'_n$ to $C'_{n+1}$,
the difference of two lifts of a fixed map to the smooth curve $C$
is a global section of
\[
                       h^*T_C\otimes(N')^{-n}.
\]
For $1\leq n<e$, its degree is
\[
               m\left(-(2g-2)+\frac{nd}{e}\right)<0
\]
by (PN1) and $d<2g-2$.  The section space is zero.  Induction proves
the claim, including the vacuous case $e=1$.

Now restrict $\phi$ to $C'_{e+1}$.  The target map factors through
the first neighborhood $C_2$ because $2e\geq e+1$.  Modulo the
$e$th power on the source, its map to $C$ is exactly the retracted
map established above.  The leading normal coefficient in (PN1)
identifies $h^*N^{-1}$ with $(N')^{-e}$.
Choose local splittings of the square-zero extension defining $C_2$.
Their differences represent $e_Y$.  Pulling these splittings back
along $h$, and comparing them with the local expressions of $\phi$
modulo the $(e+1)$st power, shows that the resulting differences are
the coboundary of the order-$e$ tangential corrections.  Hence
\[
 h^*e_Y=0
       \quad\hbox{in }H^1\bigl(C',h^*(T_C\otimes N^{-1})\bigr).
                                                               \tag{PN2}
\]
Equivalently, (PN2) is the obstruction to lifting the already retracted
map $C'_e\to C$ through $C_2$ with its prescribed leading normal map.
The possible lower tangential corrections have been excluded above;
we do not assume their vanishing without the negative-degree argument.
Since $T_C\otimes N^{-1}$ has degree $d-(2g-2)<0$, (PN2) contradicts
the curve-theoretic input.

For the final sufficient condition, set $L=T_C\otimes N^{-1}$ and
$a=-\deg L=2g-2-d>0$.  On a smooth curve over a perfect field the
differential has kernel $\mathcal O_C^p$: this is checked in the
function field using a separating parameter, and regularity follows
by normality.  Consequently there is an injection
\[
           F_*\mathcal O_C/\mathcal O_C\hookrightarrow F_*\Omega_C^1.
\]
Tensoring by $L$ and using the projection formula shows that the
kernel of
\[
                    H^1(C,L)\longrightarrow H^1(C,L^{\otimes p})
\]
is a quotient of a subspace of
$H^0(C,\Omega_C^1\otimes L^{\otimes p})$.
This last space vanishes if $pa>2g-2$.  The same inequality remains
true after replacing $L$ by any $L^{\otimes p^r}$.
All successive Frobenius pullbacks are therefore injective, so a
nonzero $e_Y$ survives every one of them.
\end{proof}

\begin{remark}
We have proved a conditional geometric obstruction and have not
chosen a numerical curve or a contracted algebraic germ.  No finite
example is introduced.  A failure of this particular cone-cover
construction does not imply that its wild inertia is infinite or that
its full fundamental group requires infinitely many generators.
\end{remark}
